\documentclass[10pt,a4paper]{amsart}
\usepackage[margin=19mm]{geometry}
\usepackage{amsmath,amssymb,mathtools}
\usepackage[scr=rsfs,cal=euler]{mathalfa}
\usepackage{xfrac}
\usepackage[unicode,hidelinks,bookmarks=false]{hyperref}
\newcommand{\ie}{i.e.\ }
\newcommand{\cf}{cf.\ }
\newcommand{\resp}{respectively\ }
\newcommand{\Subsection}{\subsection}
\providecommand{\nobreakdash}{\nobreak\hbox{-}}
\setlength{\emergencystretch}{2em}
\multlinegap0pt
\usepackage{bm}
\usepackage{bbm}
\usepackage{dsfont}
\usepackage{xspace}
\usepackage{cancel}
\usepackage{mathtools}
\usepackage{amsthm}
\makeatletter
\@ifundefined{assumption}{\newtheorem{assumption}{Assumption}}{}
\@ifundefined{lemma}{\newtheorem{lemma}{Lemma}}{}
\makeatother
\newcommand{\R}{\ensuremath{\mathbb{R}}\xspace}
\newcommand{\V}{\ensuremath{\mathcal{V}}\xspace}
\newcommand{\X}{\ensuremath{\mathcal{X}}\xspace}
\title[A Distributed Asynchronous Generalized Momentum Algorithm Wi]{A Distributed Asynchronous Generalized Momentum Algorithm Without Delay Bounds}
\author{Ellie Pond, Yichen Zhao, Matthew Hale}
\date{Source: arXiv:2508.08218v2 (CC BY 4.0)}
\begin{document}
\maketitle

\noindent\emph{Source and licence.} A Distributed Asynchronous Generalized Momentum Algorithm Without Delay Bounds. Version arXiv:2508.08218v2 (CC BY 4.0), distributed under the Creative Commons Attribution 4.0 International licence, as recorded in the arXiv licence metadata for this exact version and shown on the abstract page https://arxiv.org/abs/2508.08218v2. Full rights notice: licenses/arxiv-2508.08218-CC-BY-4.0.txt.\par\smallskip

\noindent\textit{Editorial scope.} Counted unit: the Lemma whose source label is \texttt{lem:alphabound} in the article (source0.tex lines 1272--1285, source label \texttt{lem:alphabound}), delivered together with the article's own complete original proof of it (source0.tex lines 1286--1400, one \texttt{proof} environment of 115 lines). No mathematics is written, repaired or re-derived: statement and proof are the article's own text line for line. Required context retained uncounted: eq:prob1 (lines 259--261); ass:X (lines 289--291); ass:H (lines 306--309). Every definition, display and label of those lines is retained unchanged. Omitted from this package and not redistributed: 1--258, 262--288, 292--305, 310--1271, 1401--1942. Nothing inside a retained excerpt is omitted or altered: every retained body is delivered exactly as the article's lines are written, so no omission receipt of any kind accompanies this package. Citations: the retained excerpts carry no citation command at all, so no bibliography entry of the article is transcribed and no reference is redistributed. Reference adaptation: none was needed and none was made; every reference inside the delivered unit resolves inside it, so no unresolved reference or citation is produced. Printed slips kept literal: no printed word, symbol, hypothesis or number of the counted statement or of its proof was altered. Layout and class: the article's own class is \texttt{ieeecolor} with packages including generic, cite, amsmath, amssymb, amsfonts, graphicx, xspace, textcomp, mathtools, booktabs, multirow, multicol, algorithm, xfrac; the wrapper is the lane's pinned \texttt{amsart} editorial wrapper at 10pt on A4, which restates the theorem-like environments the retained text needs and adds the article's own macro definitions that the retained text needs (\texttt{\textbackslash R}, \texttt{\textbackslash V}, \texttt{\textbackslash X}), with extra wrapper packages (bm, bbm, dsfont, xspace, cancel, mathtools). The delivered numbering is the unit's own. Assets: no retained span contains a figure environment, a graphics-inclusion command, a PDF-page inclusion, a table or a TikZ picture; no image file is copied into this package, the package declares no asset dependency and no image file is redistributed with it. Licence: version arXiv:2508.08218v2 of this article is distributed under the Creative Commons Attribution 4.0 International licence, as recorded in the arXiv licence metadata for this exact version and shown on the abstract page https://arxiv.org/abs/2508.08218v2. A full rights notice is delivered at licenses/arxiv-2508.08218-CC-BY-4.0.txt. Characters: every retained excerpt is pure ASCII, so the delivered engine, XeLaTeX with its default Latin Modern text font, reports no missing character.
    \begin{equation} \label{eq:prob1}
        \underset{x \in \X}{\operatorname{minimize}} \,\, f(x) = \sum_{i=1}^nf_i(x_{\V_i}). 
    \end{equation}    

\begin{assumption}\label{ass:X}
    The constraint set $\X\subset \R^n$ is nonempty, convex, and compact. The set can be decomposed as $\X=\X_1 \times \X_2 \times \dots \times \X_n$, where $\X_i\subset \R^{n_i}$ for all $i\in \V$.
\end{assumption}

\begin{assumption} \label{ass:H}
    The Hessian matrix~$H(x):=\nabla^2f(x)\in \R^{n\times n}$ is $\mu-$diagonally dominant on $\X\subset \R^n$ for some $\mu>0$. 
    That is, for each $i\in \V$ we have $H_{ii}(x)\geq \mu + \sum_{j=1,j\not=i}^n\vert H_{ij}(x) \vert$ for all $x\in \X$.
\end{assumption}

    \begin{lemma}\label{lem:alphabound}
        Consider the problem in~\eqref{eq:prob1} and let Assumptions \ref{ass:X}-\ref{ass:H} hold.
        Define
        \begin{align}\label{eq:alpha1}
        \alpha_1&= \big(1+\beta-\gamma\mu(1+\lambda)\big)^2 
        \\ &\quad\quad\quad + (\beta-\gamma\lambda \mu)\big(2+\beta-\gamma\mu(1+\lambda)\big)
    \end{align}
    and
    \begin{align}\label{eq:alpha2}
        \alpha_2 = 1-\gamma\mu+2(\beta-\lambda\gamma\mu).
    \end{align}
        If $(\gamma,\lambda,\beta)\in C_1\cup C_2$, then $\alpha=\max\{\alpha_1,\alpha_2\}$ 
        satisfies $\alpha\in(0,1)$.
    \end{lemma}

    \begin{proof}
        Since $(\gamma,\lambda,\beta)\in C_1\cup C_2$, we consider two cases.
    First, if $(\gamma,\lambda,\beta)\in C_1$, then
    \begin{align}
        \alpha_2 &= 1 + 2\beta - \gamma\mu(1+2\lambda)
        \\ &\leq 1 - \gamma\mu + 2\lambda(1-\gamma\mu)
        \\ &< 1 - \gamma\mu + \frac{\gamma\mu}{1-\gamma\mu}(1-\gamma\mu)
        \\ &= 1 -\gamma\mu + \gamma\mu =1,
    \end{align}
    where the first inequality is obtained by replacing $\beta$ with its upper bound $\lambda$ and the strict inequality is obtained from replacing $\lambda$ with its upper bound $\frac{\gamma\mu}{2(1-\gamma\mu)}$.
    Then
    \begin{align}
        \alpha_1 &= \big(1+\beta-\gamma\mu(1+\lambda)\big)^2
        \\ &\quad\quad\quad+(\beta-\gamma\lambda\mu)\big(2+\beta-\gamma\mu(1+\lambda)\big)
        \\&\leq \big(1 -\gamma\mu+\lambda(1-\gamma\mu)\big)^2
        \\ &\quad\quad\quad + \lambda(1-\gamma\mu)\big(2-\gamma\mu + \lambda(1-\gamma\mu)\big),
        \\ &<\Big(1-\gamma\mu + \frac{\gamma\mu(1-\gamma\mu)}{2(1-\gamma\mu)}\Big)^2
        \\ &\quad\quad\quad + \frac{\gamma\mu(1-\gamma\mu)}{2(1-\gamma\mu)}\Big(2 - \gamma\mu + \frac{\gamma\mu(1-\gamma\mu)}{2(1-\gamma\mu)}\Big)
        \\ &= \big(1-\gamma\mu + \frac{\gamma\mu}{2}\big)^2 + \frac{\gamma\mu}{2}\big(2-\gamma\mu +\frac{\gamma\mu}{2}\big)
        \\ &= \big(1-\frac{\gamma\mu}{2}\big)^2+ \frac{\gamma\mu}{2}\big(2-\frac{\gamma\mu}{2}\big)
        \\ &= 1 + \frac{(\gamma\mu)^2}{4}- \gamma\mu + \gamma\mu - \frac{(\gamma\mu)^2}{4} = 1,
    \end{align}
    where the first inequality is obtained by replacing $\beta$ with its upper bound $\lambda$ and the strict inequality comes from replacing $\lambda$ with its upper bound $\frac{\gamma\mu}{2(1-\gamma\mu)}$.
        
    Now consider the lower bound of $\alpha_2$ if $(\gamma,\lambda,\beta)\in C_1$, namely
    \begin{align}
        \alpha_2 &= 1 + 2\beta-\gamma\mu(1+2\lambda)
        \\&= 1 -\gamma\mu + 2(\beta-\gamma\lambda\mu)
        \\ &\geq 1-\gamma\mu,
    \end{align}
    where the inequality comes from multiplying the upper bound $\gamma< \frac{\beta}{\lambda\max\limits_{i\in\V}\max\limits_{\eta\in\X}\vert H_{ii}(\eta)\vert}$ by $\mu$ 
    to get $\gamma\mu< \frac{\beta}{\lambda}$ and then rearranging to get $\beta-\gamma\lambda\mu>0$.
    Using~$\beta \leq \lambda$ we see that~$\gamma\mu < \frac{\beta}{\lambda} \leq 1$ and 
    %Further, by replacing $\beta$ with its upper bound $\lambda$ in $\gamma\mu < \frac{\beta}{\lambda}$ results in $\gamma\mu<1$. Then, we obtain
    \begin{align}
        \alpha_2 &= 1-\gamma\mu>0.
    \end{align}
    Next, if $(\gamma,\lambda,\beta)\in C_1$, then we find the lower bound
    \begin{align}
        \alpha_1 &=\big(1 + \beta -\gamma\mu(1+\lambda)\big)^2 
        \\ &\quad\quad\quad + (\beta - \gamma\lambda\mu)\big(2+\beta-\gamma\mu(1+\lambda)\big)
        \\ &=\big((1-\gamma\mu) + (\beta-\gamma\lambda\mu)\big)^2
        \\ &\quad\quad\quad + (\beta-\gamma\lambda\mu)\big((2 -\gamma\mu) + (\beta-\gamma\lambda\mu)\big)^2
        >0,
    \end{align}
    where the inequalities~$1 - \gamma\mu > 0$, $2 - \gamma\mu > 0$, and
    $\beta - \gamma\lambda\mu > 0$ imply that the whole expression is strictly positive. 
    %\ep{where we replace $1-\gamma\mu$ (as well as $2-\gamma\mu$) with its strict lower bound $0$ and $\beta-\gamma\lambda\mu$ with its strict lower bound $0$.
    Then $\alpha=\max\{\alpha_1,\alpha_2\}\in (0,1)$ holds for $(\gamma,\lambda,\beta)\in C_1$.
      
    Next, if $(\gamma,\lambda,\beta)\in C_2$, then
    \begin{align}
        \alpha_2 &= 1+2\beta-\gamma\mu(1+2\lambda)
        \\ &< 1 +\gamma\mu(1+2\lambda) - \gamma\mu(1+2\lambda)=1,
    \end{align}
    where we have replaced $\beta$ with its strict upper bound $\frac{1}{2}\gamma\mu(1+2\lambda)$.
    For $\alpha_1$, we have the upper bound
    \begin{align}
        \alpha_1 &= \big(1+\beta-\gamma\mu(1+\lambda)\big)^2 
        \\ &\quad\quad\quad +  (\beta-\gamma\lambda\mu)\big(2+\beta-\gamma\mu(1+\lambda)\big)
        \\ &< \big(1 +\frac{1}{2}\gamma\mu+\lambda \gamma\mu -\gamma\mu-\lambda\gamma\mu\big)^2
        \\ &\quad\quad\quad + \big( \frac{1}{2}\gamma\mu+\lambda\gamma\mu-\lambda\gamma\mu\big)
        \\ & \quad\quad\quad \cdot\big(2 + \frac{1}{2}\gamma\mu+\lambda \gamma\mu - \gamma\mu-\lambda\gamma\mu\big)
        \\ &= \big( 1- \frac{1}{2}\gamma\mu\big)^2+ \frac{1}{2}\gamma\mu\big(2-\frac{1}{2}\gamma\mu\big)
        \\ &= 1 + \frac{1}{4}(\gamma\mu)^2 -\gamma\mu+\gamma\mu - \frac{1}{4}(\gamma\mu)^2=1,
    \end{align}
    where the inequality comes from the bound~$\beta<\frac{1}{2}\gamma\mu(1+2\lambda)$.
    We lower bound $\alpha_2$ as    
    % \begin{align}
    %     \alpha_2 &= 1 + 2\beta-\gamma\mu(1+2\lambda)
    %     \\ &> 1-\gamma\mu + 2(\beta-\lambda)
    %     \\ &\geq 1-\gamma\mu +2 (\lambda-\lambda)
    %     \\ &= 1-\gamma\mu
    % \end{align}
    \begin{align}
        \alpha_2 &= 1 + 2\beta-\gamma\mu(1+2\lambda)
            \\ &> 1+ 2\beta-(1+2\lambda)
            \\ &= 2(\beta-\lambda)\geq0,
    \end{align}
    % \ep{where the first inequality comes from~$\gamma\mu < 1$ and 
    % the second inequality is found by replacing $\beta$ with its lower bound $\lambda$.
    % We multiply the bound $\gamma<\frac{1}{\max\limits_{i\in\V}\max\limits_{\eta\in\X}\vert H_{ii}(\eta)\vert}$ 
    % by~$\mu$ to get~$\gamma\mu<1$, i.e., $1-\gamma\mu > 0$.
    % Then}
    % \begin{align}
    %     \alpha_2\geq 1-\gamma\mu>0.
    % \end{align}
    where the first inequality is found by replacing~$\gamma\mu$ with its upper bound $1$ and 
    the second inequality is due to $\beta\geq\lambda$.
    % \mh{Sure that seems fine.}
    % \yz{I think we can shorten this part by doing this:}
    % \yz{
    %     \begin{align}
    %         \alpha_2 &= 1 + 2\beta-\gamma\mu(1+2\lambda)
    %         \\ &> 1+ 2\beta-(1+2\lambda)
    %         \\ &= 2(\beta-\lambda)\geq0,
    %     \end{align}
    % where the first inequality is found by replacing~$\gamma\mu$ with its upper bound $1$ and 
    % the second inequality is due to $\beta\geq\lambda$.}
    % \mh{Sure that seems fine.}
    % \ep{Okay I'm with  you, making that change.}
    For $\alpha_1$ we have
    \begin{align}
        \alpha_1 &=\big(1+\beta-\gamma\mu(1+\lambda)\big)^2
        \\ &\quad\quad\quad + (\beta-\gamma\lambda\mu)\big(2+\beta-\gamma\mu(1+\lambda)\big)
        \\ &\geq \big(1 -\gamma\mu + \lambda(1-\gamma\mu)\big)^2
        \\ &\quad\quad\quad + \lambda(1-\gamma\mu)\big(2-\gamma\mu+\lambda(1-\gamma\mu)\big)
        \\ &\geq (1-\gamma\mu)^2 >0,
    \end{align}
    where the first inequality is found by replacing $\beta$ with its lower bound $\lambda$. 
    The second inequality is from replacing $\lambda$ with its lower bound $0$.
    The strict inequality is due to $1-\gamma\mu>0$.
    Then $\alpha=\max\{\alpha_1,\alpha_2\}\in(0,1)$ holds for $(\gamma,\lambda,\beta)\in C_2$
    as well. 
    \end{proof}

\end{document}
